\chapter{Marchés de l'inflation}\label{ch:m2:inflation-markets}

Sur les douze mois précédant juin 2022, les prix à la consommation américains ont progressé de 9.1~\%, le
rythme le plus rapide depuis plus de quarante ans. Cela aurait dû être l'année de
obligation protégée contre l'inflation, dont le principal croît avec l'indice des prix. Un
investisseur qui a acheté, le premier jour de bourse de 2022, une obligation de ce type à dix ans
et l'a évaluée le dernier jour de l'année, a vu son principal augmenter de 7.7~\% et son cours baisser de 21.3~\%~:
soit une perte de 15.2~\% sur l'année. Rien ne clochait dans l'indexation. L'obligation
protège contre l'inflation par rapport à son \emph{\omterm{def:m2:inflation-markets:breakeven}{rendement réel}}, et ce rendement
est passé de $-0.97\%$ à $+1.58\%$ lorsque la Réserve fédérale a relevé ses taux. Ce
chapitre explique ce que verse une \omterm{def:m2:inflation-markets:linker}{obligation indexée} sur l'inflation, et comment son rendement se
décompose en un taux réel et le prix de marché de l'inflation, et comment l'inflation
est échangée dans les swaps, et pourquoi le calendrier des publications de prix rend le
carry d'une position indexée sur l'inflation, mois après mois, l'un des chiffres les plus prévisibles
des titres à revenu fixe.

\section{Obligations indexées}

\begin{definition}[Obligation indexée, indice de référence de l'inflation, plancher de déflation]\label{def:m2:inflation-markets:linker}
Une \emph{obligation indexée}\index{obligation indexée} (un \emph{linker}) est une obligation
dont le principal, et donc chaque coupon, est multiplié par la croissance d'un indice des prix
à la consommation depuis sa date de jouissance. Le calcul s'appuie sur
l'\emph{indice de référence de l'inflation}\index{indice de référence de l'inflation}~: l'indice des prix d'un jour donné,
tel que les règles de l'obligation le définissent à partir des valeurs mensuelles publiées. Le rapport
de l'indice de référence un jour donné à sa valeur à la date de jouissance est le
\emph{ratio d'indexation}. Un \emph{plancher de déflation}\index{plancher de déflation} est une
promesse selon laquelle le principal remboursé à l'échéance n'est jamais inférieur à la
valeur nominale initiale, quelle qu'ait été l'évolution de l'indice.
\end{definition}

Les Treasury Inflation-Protected Securities, TIPS, ont été mis en adjudication pour la
première fois par les États-Unis en janvier 1997. Ils sont indexés sur l'indice des prix à la
consommation pour l'ensemble des consommateurs urbains, non corrigé des variations saisonnières, publié chaque mois par le
Bureau of Labor Statistics~; ils versent un coupon fixe tous les six mois sur le
principal ajusté~; à l'échéance, le porteur reçoit le principal ajusté ou le principal d'origine,
le montant le plus élevé. Le coupon fixe est un taux \emph{réel}~:
un TIPS à coupon de 0.125~\% verse 0.0625~\% de son principal ajusté à chaque semestre,
de sorte que ses flux conservent leur pouvoir d'achat.

Un indice mensuel est publié environ deux semaines après le mois qu'il mesure, si bien qu'aucune
obligation ne peut être indexée sur les prix du jour. La règle des TIPS, fixée dans la
réglementation du Trésor américain, fixe l'indice de référence du premier jour d'un mois égal à
celui du troisième mois précédent, et interpole linéairement pour les
autres jours~:
\begin{equation}\label{eq:m2:inflation-markets:refcpi}
\mathrm{Ref}(d) \;=\; I_{M-3} + \frac{t-1}{D}\,\bigl(I_{M-2} - I_{M-3}\bigr),
\end{equation}
où $d$ désigne le jour $t$ du mois $M$, $D$ le nombre de jours de $M$ et $I_m$
l'indice publié du mois $m$~; l'indice de référence et le ratio d'indexation sont
tronqués à six décimales et arrondis à cinq décimales. L'exemple du règlement lui-même~:
avec janvier 1996 à 154.40 et février à 154.90, l'\omterm{def:m2:inflation-markets:linker}{indice de référence de l'inflation} du 15
avril 1996 vaut $154.40 + \tfrac{14}{30}\times 0.50 = 154.63333$.

\begin{omfigure}
\centering
\begin{tikzpicture}[font=\small,x=1.55cm]
\draw[thick,-{Stealth}] (0,0) -- (6.6,0);
\foreach \m/\l in {0/mai,1/juin,2/juil.,3/août,4/sept.,5/oct.,6/nov.} {\draw (\m,0.1) -- (\m,-0.1) node[below] {1 \l};}
\draw[omDef,very thick] (3,1.2) -- (4,1.75);
\node[omDef,align=center,font=\footnotesize] at (3.5,2.25) {indice de référence\\en août};
\fill[omDef] (3,1.2) circle (1.8pt) node[left,font=\footnotesize] {indice de mai};
\fill[omDef] (4,1.75) circle (1.8pt) node[right,font=\footnotesize] {indice de juin};
\draw[omThm,thick,-{Stealth}] (2.39,-0.9) -- (2.39,-0.08);
\node[omThm,align=center,font=\footnotesize] at (2.39,-1.35) {13 juillet~: indice de juin\\publié};
\draw[omProp,thick,decorate,decoration={brace,amplitude=4pt,mirror}] (2.39,-1.95) -- (3,-1.95);
\node[omProp,font=\footnotesize,align=center] at (2.7,-2.45) {la progression d'août est\\connue avant le début du mois};
\end{tikzpicture}

\omcaption{Le \omterm{def:m2:inflation-markets:lag}{décalage d'indexation} d'un TIPS. L'\omterm{def:m2:inflation-markets:linker}{indice de référence de l'inflation} évolue pendant
août, de l'indice de mai à celui de juin, tous deux déjà publiés, si bien que la
croissance du principal sur août est connue dès le 13 juillet, date à laquelle l'indice de juin
a été publié en 2022. Schéma indicatif.}\label{fig:m2:inflation-markets:lag}
\end{omfigure}

\begin{example}[Un ratio d'indexation]\label{ex:m2:inflation-markets:ratio}
Prenons un TIPS fictif à dix ans, de date de jouissance le 15 janvier 2022, avec un
coupon de 0.125~\%. Son \omterm{def:m2:inflation-markets:linker}{indice de référence de l'inflation} ce jour-là
interpole entre l'indice d'octobre 2021, 276.589, et celui de novembre, 277.948~:
$276.589 + \tfrac{14}{31}\times 1.359 = 277.20274$. Le 15 août 2022, il
interpole entre mai 2022, 292.296, et juin, 296.311~: 294.10922. Le ratio
d'indexation est $294.10922/277.20274 = 1.06099$, donc 100~millions de dollars de valeur nominale
deviennent alors 106.099~millions de dollars de principal, et le coupon versé le 15 juillet, à un
ratio d'indexation de 1.04814, s'élevait à 65\,509~dollars.
\end{example}

\begin{remark}[Lorsque l'indice n'est pas publié]\label{rem:m2:inflation-markets:contingency}
Une \omterm{def:m2:inflation-markets:linker}{obligation indexée} sur une statistique dépend de la production de cette statistique. La
règle du Trésor prévoit le cas de son absence~: si l'indice d'un mois n'est pas publié à
la fin du mois suivant, le Trésor annonce un indice de remplacement, égal à la dernière valeur publiée,
prolongée d'un mois au taux annuel observé le plus récent,
$I_M = I_{M-1}\,(I_{M-1}/I_{M-13})^{1/12}$, et le conserve même si la valeur
réelle apparaît plus tard. Cela s'est produit pour la première fois en 2025~: le Bureau of
Labor Statistics n'a relevé aucun prix du 1\up{er} octobre au 12 novembre, pendant
l'interruption des crédits budgétaires fédéraux, et aucun indice d'octobre n'a été publié. Le 26
novembre, le Trésor a annoncé 325.604 pour octobre, ce qui correspond à
$324.800 \times (324.800/315.301)^{1/12}$ à partir des indices de septembre 2025
et de 2024. Chaque flux de TIPS qui dépend d'octobre 2025 utilise ce nombre.
\end{remark}

Les autres grands marchés reposent sur la même conception, avec d'autres indices. Au
Royaume-Uni, les gilts indexés sont liés au Retail Prices Index~; ceux émis avant 2005
utilisent un décalage de huit mois~; ceux émis depuis ont un décalage de trois mois
assorti d'un indice de référence interpolé, et un gilt rembourse moins que sa valeur nominale
si l'indice a baissé~: il n'existe pas de \omterm{def:m2:inflation-markets:linker}{plancher de déflation}.
La France émet deux familles, l'une indexée sur les prix à la consommation française hors
tabac, et l'autre, depuis 2001, sur l'indice harmonisé de la zone euro hors tabac, avec la
même interpolation entre les indices publiés il y a trois mois
et il y a deux mois.

\begin{dated}{2026-09}{Marchés des obligations indexées}\label{dat:m2:inflation-markets:markets}
\textbf{États-Unis}~: TIPS négociables en circulation pour 2\,152.7~milliards de dollars le 31
août 2026, soit 6.8~\% des 31\,828.0~milliards de dollars de dette négociable du Trésor~; échéances
de 5, 10 et 30 ans. \textbf{Royaume-Uni}~: le Retail Prices Index doit être calculé selon les
méthodes et les données de l'indice des prix à la consommation, y compris les coûts du logement
des propriétaires occupants (CPIH) à partir de février 2030, comme l'ont annoncé la
Statistics Authority et le Trésor le 25 novembre 2020, une date choisie pour
limiter l'effet sur les détenteurs de gilts indexés. \textbf{Zone euro}~:
Eurostat a rebasé l'indice harmonisé sur l'année 2025 à partir de février 2026~; le Trésor
français indique que ce changement de base ne modifie pas les coefficients d'indexation de ses obligations.
\end{dated}

\section{Rendements réels et points morts}

\begin{definition}[Rendement réel, point mort d'inflation]\label{def:m2:inflation-markets:breakeven}
Le \emph{rendement réel}\index{rendement réel} d'une \omterm{def:m2:inflation-markets:linker}{obligation indexée} est le rendement
actuariel de ses flux réels (le coupon fixe et la valeur faciale, exprimés en unités
de l'indice du jour) à son prix réel, c'est-à-dire le prix coté pour 100 du
principal ajusté. Le \emph{point mort d'inflation}\index{point mort d'inflation}
à une échéance donnée est le taux d'inflation constant pour lequel une obligation nominale et un titre indexé
de même échéance rapportent le même rendement~:
\begin{equation}\label{eq:m2:inflation-markets:fisher}
1 + y_{\mathrm{nom}} \;=\; (1 + y_{\mathrm{real}})(1 + \pi_{\mathrm{BE}}),
\qquad \pi_{\mathrm{BE}} \approx y_{\mathrm{nom}} - y_{\mathrm{real}}.
\end{equation}
\end{definition}

Le 29 juillet 2022, le rendement nominal du Trésor à dix ans s'élevait à 2.67~\% et le \omterm{def:m2:inflation-markets:breakeven}{rendement
réel} à dix ans à 0.14~\%~: le point mort exact vaut $1.0267/1.0014 - 1 = 2.5265\%$ et l'écart
2.53~\%, comme l'affichent les écrans de données. Un investisseur qui anticipe une inflation
supérieure au point mort préfère le titre indexé~; celui qui anticipe une inflation plus faible préfère l'obligation nominale.

\begin{proposition}[Ce que contient un point mort]\label{prop:m2:inflation-markets:contents}
Sous la mesure neutre au risque, le point mort n'est égal à l'inflation anticipée à cet horizon
que si les investisseurs n'exigent aucune prime pour porter le risque d'inflation et si les
deux obligations sont aussi liquides l'une que l'autre. En général
\[
\pi_{\mathrm{BE}} \;\approx\; \mathbb{E}[\pi] + \text{prime de risque d'inflation}
- \text{prime de liquidité} + \text{termes de convexité},
\]
avec un terme de liquidité qui entre négativement, car un titre indexé moins liquide doit
offrir un \omterm{def:m2:inflation-markets:breakeven}{rendement réel} plus élevé, ce qui abaisse le point mort mesuré.
\end{proposition}
\admitted{La décomposition est un énoncé sur le noyau de tarification~; les modèles de structure par terme
qui estiment chacune de ses parties font l'objet du livre~6.}

La conséquence pratique~: un point mort est un prix, non une prévision. Il recule
lorsque les investisseurs se réfugient vers les actifs les plus liquides, comme en mars 2020,
où le point mort à dix ans est tombé à 0.50~\% le 19 mars, et il peut rester
inférieur à l'inflation réalisée pendant de longues périodes.

\begin{omfigure}
\begin{tikzpicture}
\begin{axis}[width=11cm,height=5.2cm,
  xlabel={},ylabel={pourcent},
  tick label style={font=\small},label style={font=\small},
  xmin=2018,xmax=2026.75,ymin=-1.5,ymax=5.5,grid=major,grid style={gray!25},
  xtick={2018,2019,2020,2021,2022,2023,2024,2025,2026},xticklabel style={/pgf/number format/1000 sep={}},
  legend columns=3,legend style={font=\footnotesize,at={(0.5,-0.18)},anchor=north,draw=none,fill=none,/tikz/every even column/.append style={column sep=8pt}},
  table/col sep=comma]
\addplot[omDef,very thick] table[x=t,y=nominal]{figdata/markets-2/11-inflation-markets/tips10.csv};
\addplot[omThm,very thick] table[x=t,y=real]{figdata/markets-2/11-inflation-markets/tips10.csv};
\addplot[omProp,thick,dashed] table[x=t,y=breakeven]{figdata/markets-2/11-inflation-markets/tips10.csv};
\draw[gray] (axis cs:2018,0) -- (axis cs:2026.75,0);
\legend{nominal à dix ans,réel à dix ans (TIPS),point mort à dix ans}
\end{axis}
\end{tikzpicture}

\omcaption{Rendements nominal et réel du Trésor à dix ans et point mort qui les sépare,
fins de mois, de janvier 2018 à août 2026. En 2022, le \omterm{def:m2:inflation-markets:breakeven}{rendement réel} a progressé de
plus de deux points et demi de pourcentage, alors que le point mort a à peine bougé~:
les pertes des \omterm{def:m2:inflation-markets:linker}{obligations indexées} sur l'inflation cette année-là proviennent des taux réels. Données~: FRED
séries DGS10, DFII10 et T10YIE.}\label{fig:m2:inflation-markets:tips10}
\end{omfigure}

La figure répond aussi à l'accroche. Du début à la fin de 2022, le
point mort a légèrement reculé, de 2.60~\% à 2.30~\%, tandis que le \omterm{def:m2:inflation-markets:breakeven}{rendement réel} progressait de 2.55
points de pourcentage. Le cours d'un titre indexé évolue avec son \omterm{def:m2:inflation-markets:breakeven}{rendement réel} exactement comme
celui d'une obligation nominale évolue avec son rendement nominal, avec une duration réelle de même ordre~;
l'indexation s'ajoute par-dessus. Le TIPS à dix ans de l'accroche avait une duration réelle
supérieure à neuf ans, si bien que ce mouvement lui a coûté environ un cinquième de sa valeur,
alors qu'une année d'inflation à 9~\%, qui atteint le principal avec un décalage, lui a ajouté 7.7~\%.

\section{Swaps d'inflation zéro-coupon}

\begin{definition}[Swap d'inflation zéro-coupon]\label{def:m2:inflation-markets:zciis}
Un \emph{swap d'inflation zéro-coupon}\index{swap d'inflation zéro-coupon} de notionnel
$N$, de maturité $T$ années et de taux fixe $K$ échange un unique paiement net à
l'échéance~: le receveur d'inflation reçoit
\[
N\left(\frac{I(T)}{I(0)} - (1+K)^T\right),
\]
où $I(0)$ et $I(T)$ sont l'indice de référence à la date de départ et à l'échéance,
avec le même décalage que les obligations. Rien n'est échangé au départ.
\end{definition}

Les swaps d'inflation en dollars font référence au même indice non désaisonnalisé que
les TIPS, et le swap zéro-coupon est la forme la plus courante~: une étude des transactions des dealers
menée par la Federal Reserve Bank of New York a dénombré 144 transactions de juin à
août 2010, soit environ 65~millions de dollars par jour, contre environ 5~milliards de dollars par jour
pour les TIPS. Les fonds de pension et les assureurs, dont les engagements croissent avec les prix, reçoivent
l'inflation~; les émetteurs dont les revenus sont indexés, comme les fournisseurs d'énergie, la paient.

\begin{proposition}[L'indice à terme]\label{prop:m2:inflation-markets:forward}
Le swap valant zéro à l'origine, le \omterm{def:m2:government-bonds:zero}{taux zéro-coupon} de marché $K_T$
fixe une valeur à terme de l'indice~: $I^{\mathrm{fwd}}(T) = I(0)(1+K_T)^T$ est le
niveau auquel la jambe d'inflation et la jambe fixe sont égales, et tout
flux indexé $c\,I(T)/I(0)$ payé en $T$ vaut
$c(1+K_T)^T P(0,T)$ aujourd'hui, où $P(0,T)$ est le facteur d'actualisation nominal.
\end{proposition}

\begin{proof}
Recevoir le flux indexé et payer l'inflation sur un swap de notionnel $c$~:
les termes d'indice se compensent et il reste $c(1+K_T)^T$, connu aujourd'hui, à la date $T$. Sa
valeur est $c(1+K_T)^T P(0,T)$~; le swap était gratuit.
\end{proof}

\begin{example}[Un swap à cinq ans]\label{ex:m2:inflation-markets:swap}
Un fonds de pension reçoit l'inflation sur 100~millions de dollars pendant cinq ans à 2.50~\%,
avec un indice de base de 300.00. L'indice à terme vaut $300 \times 1.025^5 = 339.42$.
Si l'indice termine à 345.00, soit 15~\% au-dessus de sa base, le fonds reçoit
$100 \times (1.15 - 1.025^5) = \text{USD}~1.86$~millions de dollars~; s'il avait
terminé au niveau à terme, rien.
\end{example}

La proposition transforme une courbe de taux de swap en une courbe de valeurs d'indice
à terme, donc en un prix pour chaque \omterm{def:m2:inflation-markets:linker}{obligation indexée}~: il suffit d'actualiser ses flux indexés
à la courbe nominale. L'écart entre ce prix et le cours de marché, coté
comme un spread par rapport à la courbe de swap, varie avec le coût, pour le bilan, de la détention
d'obligations, comme le fait le \omterm{def:m2:interest-rate-swaps:spread}{spread de swap} du \cref{ch:m2:interest-rate-swaps}
(le \cref{ch:m2:corporate-bonds} revient sur de tels spreads).

\section{Décalage d'indexation et saisonnalité}

\begin{definition}[Décalage d'indexation, saisonnalité de l'inflation]\label{def:m2:inflation-markets:lag}
Le \emph{décalage d'indexation}\index{décalage d'indexation} d'une \omterm{def:m2:inflation-markets:linker}{obligation indexée} est le délai entre
le mois dont les prix fixent son indice de référence et le mois auquel cette
référence s'applique~: trois mois pour les TIPS et pour la plupart des obligations émises depuis
les années 2000. La \emph{saisonnalité de l'inflation}\index{saisonnalité de l'inflation} est la configuration récurrente
d'un indice des prix non désaisonnalisé au cours de l'année civile, qui résulte du calendrier
des soldes, de la consommation d'énergie et des revalorisations annuelles des prix.
\end{definition}

Deux conséquences en découlent, et toutes deux comptent pour un trader. La première figure dans la
\cref{fig:m2:inflation-markets:lag}~: la croissance de l'indice de référence sur un mois est celle
de l'indice publié deux et trois mois plus tôt,
\begin{equation}\label{eq:m2:inflation-markets:accrual}
\frac{\mathrm{Ref}(\text{1er de } M+1)}{\mathrm{Ref}(\text{1er de } M)} - 1 \;=\; \frac{I_{M-2}}{I_{M-3}} - 1,
\end{equation}
si bien que l'accumulation d'inflation d'une \omterm{def:m2:inflation-markets:linker}{obligation indexée} pour le mois suivant, voire les deux suivants, est connue exactement. En
2022, les mois aux plus fortes publications se sont enchaînés~: l'indice de référence
a progressé de 1.37~\% sur le seul mois d'août, soit 17.8~\% en rythme annuel, puis a légèrement
reculé en septembre (\cref{fig:m2:inflation-markets:accrual}).

\begin{omfigure}
\begin{tikzpicture}
\begin{axis}[width=11cm,height=4.8cm,ybar,bar width=4pt,
  ylabel={accumulation sur le mois (\%)},
  tick label style={font=\small},label style={font=\small},
  xmin=-1,xmax=36,ymin=-0.5,ymax=1.95,ytick={-0.5,0,0.5,1,1.5},grid=major,grid style={gray!25},
  xtick={0,12,24},xticklabels={janv. 2021,janv. 2022,janv. 2023},
  table/col sep=comma]
\addplot[fill=omDef!70,draw=omDef] table[x=k,y=accrual]{figdata/markets-2/11-inflation-markets/accrual.csv};
\draw[omThm,thick,-{Stealth}] (axis cs:17.3,1.72) -- (axis cs:18.8,1.45);
\node[omThm,font=\footnotesize,anchor=east,align=right] at (axis cs:17.3,1.66) {août 2022~: 1.37~\%,\\connu le 13 juillet};
\end{axis}
\end{tikzpicture}

\omcaption{Croissance de l'indice de référence des TIPS sur chaque mois civil, de janvier
2021 à décembre 2023~: la variation mensuelle de l'indice non désaisonnalisé
de trois et deux mois plus tôt. Données~: série FRED CPIAUCNS, tutoriel du
chapitre.}\label{fig:m2:inflation-markets:accrual}
\end{omfigure}

La seconde conséquence tient à la saisonnalité. Comme l'indice n'est pas
désaisonnalisé, l'accumulation mensuelle présente une configuration calendaire qui se superpose à la tendance~:
les prix ont tendance à monter plus vite en début d'année et plus lentement, voire à baisser, en fin d'année
(\cref{fig:m2:inflation-markets:seasonal}). Sur une obligation à dix ans, la configuration
s'annule en moyenne. Sur quelques mois, elle ne s'annule pas, et un \omterm{def:m2:inflation-markets:breakeven}{point mort d'inflation} à court terme
lu sans correction donne un mauvais signal~: une obligation dont l'accumulation restante
couvre les mois de référence d'octobre à décembre porte environ 0.85~point de pourcentage
d'inflation de moins que le trimestre moyen, soit 3.4 points en rythme annuel.

\begin{omfigure}
\begin{tikzpicture}
\begin{axis}[width=11cm,height=4.6cm,ybar,bar width=12pt,
  ylabel={écart de variation (pp)},xlabel={mois de l'indice},
  tick label style={font=\small},label style={font=\small},
  xmin=0.3,xmax=12.7,ymin=-0.45,ymax=0.4,grid=major,grid style={gray!25},
  xtick={1,2,3,4,5,6,7,8,9,10,11,12},xticklabels={janv.,févr.,mars,avr.,mai,juin,juil.,août,sept.,oct.,nov.,déc.},
  table/col sep=comma]
\addplot[fill=omThm!60,draw=omThm] table[x=month,y=factor]{figdata/markets-2/11-inflation-markets/seasonal.csv};
\end{axis}
\end{tikzpicture}

\omcaption{Saisonnalité de l'indice des prix à la consommation américain, non désaisonnalisé~:
la variation moyenne de l'indice de chaque mois par rapport au mois précédent, moins la
variation mensuelle moyenne de son année, de 2010 à 2024 (2025 ne dispose pas de son
indice d'octobre), en points de pourcentage. Mars
est le mois le plus fort, novembre et décembre les plus faibles. Données~: série FRED CPIAUCNS,
tutoriel du chapitre.}\label{fig:m2:inflation-markets:seasonal}
\end{omfigure}

\begin{proposition}[Carry d'une obligation indexée financée]\label{prop:m2:inflation-markets:carry}
Une \omterm{def:m2:inflation-markets:linker}{obligation indexée} de valeur de marché $V$ détenue pendant $\tau$ jours, à \omterm{def:m2:inflation-markets:breakeven}{rendement réel}
inchangé $y$, financée en repo au taux $r$, rapporte
\[
\mathrm{carry} \;=\; V\left[(1+a)(1+y)^{\tau/365} - 1\right] - V r\,\frac{\tau}{360},
\]
où $a$ est la croissance de l'indice de référence sur la période, connue à l'avance
pour deux mois au plus.
\end{proposition}

\begin{proof}
À \omterm{def:m2:inflation-markets:breakeven}{rendement réel} inchangé, le prix réel croît au rythme du \omterm{def:m2:inflation-markets:breakeven}{rendement réel} (le
rappel vers le pair plus l'accumulation du coupon), et le \omterm{def:m2:bond-futures:cf}{prix de facturation} est le prix réel
multiplié par le ratio d'indexation, qui croît de $1+a$. Les intérêts de la repo sont versés sur $V$
en base actual/360.
\end{proof}

Comme $a$ est connu, il entre dans le prix~: une vente à terme de l'\omterm{def:m2:inflation-markets:linker}{obligation indexée} pour la
fin du mois est fixée à un prix qui l'intègre, et un acheteur qui précède une
publication importante le paie. Ce qui reste incertain, et ce qu'un trader gagne ou
perd, c'est le mouvement du \omterm{def:m2:inflation-markets:breakeven}{rendement réel} et les publications pas encore parues.

\section{Tutoriel~: points morts, saisonnalité et carry}\label{tut:m2:inflation-markets:carry}

\begin{tutorial}
\textbf{Objectif.} Calculer les indices de référence et les ratios d'indexation selon la
règle du Trésor, estimer la configuration saisonnière de l'indice, et évaluer le carry d'un
TIPS financé. \textbf{État final~:}
\cref{fig:m2:inflation-markets:tips10,fig:m2:inflation-markets:accrual,fig:m2:inflation-markets:seasonal},
\cref{ex:m2:inflation-markets:ratio} et les chiffres du problème du week-end.
\begin{tutsteps}
\item \textbf{L'indice de référence}, le ratio d'indexation et l'accumulation du mois,
par les \cref{eq:m2:inflation-markets:refcpi,eq:m2:inflation-markets:accrual}.
\omcode{code/firm/breakeven/firm_breakeven.py}{26}{41}{Indice de référence, ratio d'indexation et accumulation mensuelle.}\label{lst:m2:inflation-markets:ref}
\item \textbf{Saisonnalité et carry.} Variation logarithmique de chaque mois moins la moyenne de son année,
moyennée sur les années complètes (une année avec un mois manquant est écartée)~; le carry de
la \cref{prop:m2:inflation-markets:carry}.
\omcode{code/firm/breakeven/firm_breakeven.py}{59}{81}{Facteurs saisonniers et carry d'une obligation indexée financée.}\label{lst:m2:inflation-markets:carry}
\item \textbf{Exécuter} \texttt{inflation\_demo.hook\_2022()},
\texttt{inflation\_demo.august\_2022()} et \texttt{fig\_inflation.py}, qui
lisent \texttt{data/markets-2/cpi\_us\_monthly.csv} et
\texttt{tips10\_monthend.csv}.
\end{tutsteps}
\textbf{Pour aller plus loin.} Estimer les facteurs saisonniers sur 2010--2019 seulement
et comparer~; puis calculer le carry de septembre 2022 et expliquer son signe.
\end{tutorial}

\section{Construction~: le calculateur de point mort et de carry}\label{bld:m2:inflation-markets:breakeven}

\begin{build}
\bfield{Rôle} La firme miniature cote et finance les \omterm{def:m2:inflation-markets:linker}{obligations indexées}
et les swaps d'inflation~: elle a besoin de l'indice de référence de n'importe quel jour, du ratio d'indexation de chaque obligation,
des points morts, des valeurs à terme de l'indice déduites des taux de swap, et le carry
qu'elle touchera au cours du mois suivant.
\bfield{Interface} \texttt{ref\_index(d, index, lag)}~;
\texttt{index\_ratio(d, dated, index, lag)}~; \texttt{monthly\_accrual(year,
month, index, lag)}~; \texttt{breakeven(nominal, real)}~;
\texttt{forward\_index(base, zc\_rate, years)}~; \texttt{zc\_swap\_payment}~;
\texttt{seasonal\_factors(index)}~; \texttt{linker\_carry(value, accrual,
real\_yield, repo, days)}~; \texttt{contingency\_index(last,
year\_before\_last, missing)}. \bfield{Règles} Valeurs d'indice repérées par (année, mois)~; indice de référence et ratio d'indexation
tronqués à six décimales puis arrondis à cinq~; un paramètre de décalage afin que
les gilts à huit mois soient évalués par le même code~; les facteurs saisonniers
ne sont calculés qu'à partir d'années civiles complètes.
\bfield{Tests d'acceptation} \texttt{code/firm/breakeven/tests/}~: l'exemple
propre à la réglementation (154.63333 et 1.00011)~; le premier jour d'un mois égale
l'indice décalé~; le substitut d'octobre 2025 (325.604)~; un swap conclu au taux à terme ne rapporte rien~; les facteurs saisonniers somment à zéro~;
les signes du carry.
\bfield{Pour aller plus loin} Courbes d'indice à terme désaisonnalisées, déduites des taux de swap~; un
solveur de \omterm{def:m2:inflation-markets:breakeven}{rendement réel} pour un TIPS à partir de son prix réel coté~; les gilts britanniques à huit
mois de décalage, dont les coupons sont fixés à l'avance.
\end{build}

\begin{omsources}
\item Code of Federal Regulations, 31 CFR Part 356, Appendix B (formules pour
titres protégés contre l'inflation)~; TreasuryDirect, TIPS and the history of TIPS.
\item UK Debt Management Office, FAQs on index-linked gilts~; UK Statistics
Authority, réponse à la consultation conjointe sur la réforme du Retail Prices
Index, 25 novembre 2020.
\item Agence France Trésor, caractéristiques de OATi et OAT\texteuro i.
\item M.~J.~Fleming and J.~R.~Sporn, ``Trading activity and price transparency in the inflation swap market'', \emph{Federal Reserve Bank of New
York Economic Policy Review}, mai 2013.
\item Bureau of Labor Statistics, Consumer Price Index, communiqué de presse pour juin
2022~; séries FRED CPIAUCNS, CPIAUCSL, DGS10, DFII10, T10YIE~; Treasury, Monthly Statement of the Public Debt.
\end{omsources}

\section{Exercices}

\begin{exercise}[$\star$]\label{exo:m2:inflation-markets:1}
Avec l'indice de mai 2022 à 292.296 et celui de juin à 296.311, donner
l'indice de référence du TIPS au 15 août 2022.
\end{exercise}

\begin{exercise}[$\star$]\label{exo:m2:inflation-markets:2}
Le rendement nominal à dix ans est de 2.67~\% et le \omterm{def:m2:inflation-markets:breakeven}{rendement réel} à dix ans de 0.14~\%. Donner
le point mort exact ainsi que son approximation.
\end{exercise}

\begin{exercise}[$\star$]\label{exo:m2:inflation-markets:3}
Sur la durée de vie d'une obligation, l'indice recule de 3~\%. Que rembourse un TIPS pour 100
de valeur nominale, et que rembourse un gilt indexé~?
\end{exercise}

\begin{exercise}[$\star\star$]\label{exo:m2:inflation-markets:4}
Un fonds reçoit l'inflation sur 100 millions de dollars pendant cinq ans à 2.50~\%, indice de
base 300.00. L'indice finit à 345.00. Quel est le paiement net, et à quel
indice final celui-ci serait-il nul~?
\end{exercise}

\begin{exercise}[$\star\star$]\label{exo:m2:inflation-markets:5}
Donner le ratio d'indexation au 15 août 2022 du TIPS de
l'\cref{ex:m2:inflation-markets:ratio}, et expliquer pourquoi le coupon versé le 15
juillet utilisait un ratio plus petit.
\end{exercise}

\begin{exercise}[$\star\star$]\label{exo:m2:inflation-markets:6}
Avec les facteurs de la \cref{fig:m2:inflation-markets:seasonal}, de combien la
saisonnalité réduit-elle l'inflation qu'un TIPS accumule sur les trois mois dont
les indices de référence sont ceux d'octobre, de novembre et de décembre, et de combien
à taux annuel~?
\end{exercise}

\begin{exercise}[$\star\star\star$]\label{exo:m2:inflation-markets:7}
\emph{Programmation.} Estimer les facteurs saisonniers sur 2010--2019 seulement. Quels sont les
mois les plus forts et les plus faibles, et de combien un mois quelconque s'écarte de
l'estimation 2010--2024~?
\end{exercise}

\begin{exercise}[$\star\star\star$]\label{exo:m2:inflation-markets:8}
\emph{Repérer l'erreur.} «~L'inflation est à 9~\% et le point mort à dix ans
est de 2.5~\%~: acheter des TIPS, vendre des \omterm{def:m2:money-markets:bill}{bons du Trésor} nominaux, et encaisser la différence.~»
Corriger.
\end{exercise}

\section{Problème~: le carry d'une obligation indexée}

\begin{problem}[{Problème du week-end --- ce que rapporte un TIPS sur un mois de publication importante}]\label{pb:m2:inflation-markets:1}
Le 13 juillet 2022, l'indice de juin a été publié à 296.311, soit 1.37~\% au-dessus de celui de mai,
292.296. Un trader détient 100 millions de dollars de valeur nominale du TIPS à dix ans de
l'\cref{ex:m2:inflation-markets:ratio} (coupon 0.125~\%, daté du 15 janvier 2022)
et le finance en repo à 2.30~\%, soit à peu près le taux au jour le jour du début août. Le
29 juillet, le \omterm{def:m2:inflation-markets:breakeven}{rendement réel} à dix ans était de 0.14~\% et le rendement nominal de 2.67~\%~;
on suppose qu'ils ne bougent pas pendant le mois d'août.

\medskip\noindent\textbf{Partie I --- L'indice.}
\begin{enumerate}
\item Donner l'indice de référence au 1\up{er} août et au 1\up{er} septembre 2022.
\item Donner le ratio d'indexation au 1\up{er} août.
\item Donner la croissance de l'indice de référence sur le mois d'août.
\item Quand le trader en a-t-il eu connaissance, et pourquoi~?
\item Donner l'indice de référence au 15 août.
\end{enumerate}

\medskip\noindent\textbf{Partie II --- Le carry.}
\begin{enumerate}[resume]
\item Donner le prix réel au taux de 0.14~\% et la valeur de marché de la position au 1\up{er} août.
\item Donner le revenu de l'accumulation d'indice sur les 31 jours d'août.
\item Donner le revenu du \omterm{def:m2:inflation-markets:breakeven}{rendement réel}.
\item Donner le coût du repo.
\item Donner le carry.
\end{enumerate}

\medskip\noindent\textbf{Partie III --- Comparaisons.}
\begin{enumerate}[resume]
\item Donner le carry de la même valeur de marché d'une obligation nominale offrant un rendement de 2.67~\%.
\item Donner le couru d'août à taux annuel, et le comparer au
\omterm{def:m2:inflation-markets:breakeven}{point mort d'inflation}.
\item L'indice de juillet valait 296.276. Donner le couru de septembre et le carry (30 jours).
\item Donner le DV01 de la position en \omterm{def:m2:inflation-markets:breakeven}{rendement réel}, et le mouvement de \omterm{def:m2:inflation-markets:breakeven}{rendement réel} qui
annulerait le carry d'août.
\item Pourquoi le carry est-il plus élevé pour un TIPS à courte échéance, rapporté à son risque~?
\end{enumerate}

\medskip\noindent\textbf{Partie IV --- Jugement.}
\begin{enumerate}[resume]
\item Le carry constitue-t-il un profit pour un trader qui achète le 29 juillet~?
\item Pourquoi les TIPS ont-ils perdu de l'argent en 2022 malgré de tels mois~?
\item Que vaut le \omterm{def:m2:inflation-markets:linker}{plancher de déflation} de cette obligation en août 2022~?
\item Énoncer le \emph{résultat nommé}~: le carry de 100~millions de dollars de valeur nominale du
TIPS à dix ans sur août 2022.
\item En une phrase~: contre quoi une \omterm{def:m2:inflation-markets:linker}{obligation indexée} protège-t-elle, et contre quoi ne protège-t-elle pas~?
\end{enumerate}
\end{problem}

\section{Questions d'entretien}

\begin{interviewq}[$\star$ \iqroles{trader, chercheur}]\label{iq:m2:inflation-markets:1}
Qu'est-ce qu'un \omterm{def:m2:inflation-markets:breakeven}{point mort d'inflation}, et pourquoi n'est-il pas la
prévision d'inflation du marché~?
\end{interviewq}

\begin{interviewq}[$\star$ \iqroles{banque, développeur}]\label{iq:m2:inflation-markets:2}
Comment le principal d'un TIPS évolue-t-il, et que se passe-t-il à
l'échéance après une déflation~?
\end{interviewq}

\begin{interviewq}[$\star\star$ \iqroles{trader}]\label{iq:m2:inflation-markets:3}
Pourquoi un TIPS porte-t-il autant de carry dans le mois qui suit une forte publication d'inflation, et
qui perçoit ce carry~?
\end{interviewq}

\begin{interviewq}[$\star\star$ \iqroles{chercheur, trader}]\label{iq:m2:inflation-markets:4}
Le taux du \omterm{def:m2:inflation-markets:zciis}{swap d'inflation zéro-coupon} à cinq ans et le \omterm{def:m2:inflation-markets:breakeven}{point mort
d'inflation} à cinq ans du TIPS diffèrent. Qu'est-ce qui pourrait les faire diverger~?
\end{interviewq}

\begin{interviewq}[$\star\star$ \iqroles{chercheur}]\label{iq:m2:inflation-markets:5}
Comment retireriez-vous la saisonnalité d'un indice des prix pour valoriser des
\omterm{def:m2:inflation-markets:linker}{obligations indexées} à courte échéance~?
\end{interviewq}

\begin{interviewq}[$\star\star\star$ \iqroles{développeur, banque}]\label{iq:m2:inflation-markets:6}
Concevez le service qui publie le ratio d'indexation de chaque \omterm{def:m2:inflation-markets:linker}{obligation indexée} qu'une firme détient,
aux États-Unis, au Royaume-Uni et dans la zone euro, chaque matin.
\end{interviewq}
